%Daten zur Institution

%\input{Dozentdaten}

%\renewcommand{\fachbereich}{Fachbereich}

%\renewcommand{\dozent}{Prof. Dr. . }

%Klausurdaten

\renewcommand{\klausurgebiet}{ }

\renewcommand{\klausurtyp}{ }

\renewcommand{\klausurdatum}{ . 20}

\klausurvorspann {\fachbereich} {\klausurdatum} {\dozent} {\klausurgebiet} {\klausurtyp}

%Daten für folgende Punktetabelle

\renewcommand{\aeins}{  3  }

\renewcommand{\azwei}{  3   }

\renewcommand{\adrei}{  0  }

\renewcommand{\avier}{  0  }

\renewcommand{\afuenf}{  4  }

\renewcommand{\asechs}{  5  }

\renewcommand{\asieben}{  4  }

\renewcommand{\aacht}{  0  }

\renewcommand{\aneun}{  4  }

\renewcommand{\azehn}{  6  }

\renewcommand{\aelf}{  2  }

\renewcommand{\azwoelf}{  6  }

\renewcommand{\adreizehn}{  4   }

\renewcommand{\avierzehn}{  10  }

\renewcommand{\afuenfzehn}{  10  }

\renewcommand{\asechzehn}{  0  }

\renewcommand{\asiebzehn}{  61  }

\renewcommand{\aachtzehn}{    }

\renewcommand{\aneunzehn}{    }

\renewcommand{\azwanzig}{    }

\renewcommand{\aeinundzwanzig}{    }

\renewcommand{\azweiundzwanzig}{    }

\renewcommand{\adreiundzwanzig}{    }

\renewcommand{\avierundzwanzig}{    }

\renewcommand{\afuenfundzwanzig}{    }

\renewcommand{\asechsundzwanzig}{    }

\punktetabellesechzehn

\klausurnote

\newpage

\setcounter{section}{K}

__NOEDITSECTION__

\inputaufgabeklausurloesung
{3}

{

Definiere die folgenden
\zusatzklammer {kursiv gedruckten} {} {} Begriffe.
\aufzaehlungsechs{Die
\stichwort {Weingartenabbildung} {}
in $T_PY$ zu einem Punkt

\mavergleichskette
{\vergleichskette
{ P
}
{ \in }{ Y
}
{  }{
}
{    }{
}
{   }{
}
}
{}{}{}
auf einer orientierten differenzierbaren Hyperfläche

\mavergleichskette
{\vergleichskette
{ Y
}
{ \subseteq }{ \R^n
}
{  }{
}
{    }{
}
{   }{
}
}
{}{}{.}

}{Eine
\stichwort {reelles Vektorbündel} {}
$V$ vom Rang $r$ über einem
\definitionsverweis {topologischen Raum}{}{}
$X$.

}{Das \stichwort {Wegintegral} {} zu einer $1$-Differentialform
\mathl{\omega \in
{ \mathcal E }^{  1 } (  M   )}{} auf einer differenzierbaren Mannigfaltigkeit $M$ bezüglich einer stetig differenzierbaren Kurve
\maabb {\gamma} {[a,b]} {M
} {.}

}{Eine
\stichwort {positive Volumenform} {}
$\omega$ auf einer differenzierbaren Mannigfaltigkeit $M$.

}{Eine \stichwort {geschlossene} {} Differentialform $\omega$ auf einer differenzierbaren Mannigfaltigkeit $M$.

}{Die
\stichwort {tangentiale Beschleunigung} {}
einer zweifach stetig differenzierbaren Kurve
\maabb {\gamma} {I} {M
} {}
auf einer
\definitionsverweis {riemannschen Mannigfaltigkeit}{}{}
$M$.
}

}
{

\aufzaehlungsechs{Es sei $N$ ein
\definitionsverweis {Einheitsnormalenfeld}{}{}
auf $Y$, das die Orientierung festlegt. Dann nennt man
\maabbeledisp {L_P} {T_P Y} { T_PY
} {v} {- \left(  D_{v} N  \right) \left(  P  \right)
} {,}
die
Weingartenabbildung
in $T_PY$.
}{Ein
reelles Vektorbündel
$V$ vom Rang $r$ ist ein topologischer Raum zusammen mit einer
\definitionsverweis {stetigen Abbildung}{}{}
\maabb {p} {V} {X
} {}
derart, dass jede
\definitionsverweis {Faser}{}{}
$p^{-1}(x)$ ein
$r$-\definitionsverweis {dimensionaler}{}{}
\definitionsverweis {reeller Vektorraum}{}{}
ist und dass es eine
\definitionsverweis {offene Überdeckung}{}{}

\mavergleichskette
{\vergleichskette
{X
}
{ = }{ \bigcup_{i \in I} U_i
}
{  }{
}
{    }{
}
{   }{
}
}
{}{}{}
und
\definitionsverweis {Homöomorphismen}{}{}
\maabbdisp {\varphi_i} {p^{-1}  \left(  U_i  \right) } { U_i \times \R^r
} {}
über $U_i$ gibt, die in jeder Faser einen
\definitionsverweis {linearen Isomorphismus}{}{}
\maabbdisp {\left(  \varphi_i  \right)_x} {p^{-1} (x) } { \R^r
} {}
induzieren.
}{Das Wegintegral ist durch

\mavergleichskettedisp
{\vergleichskette
{  \int_\gamma \omega
}
{ =} {\int_a^b \gamma^* \omega
}
{ } {
}
{ } {
}
{ } {
}
}
{}{}{}
definiert.
}{Eine
$n$-\definitionsverweis {Differentialform}{}{}
$\omega$ auf $M$ heißt eine positive Volumenform, wenn für jede
\definitionsverweis {Karte}{}{}
\maabbdisp {\alpha} { U } { V
} {}
\zusatzklammer {mit \mathlk{V \subseteq \R^n}{} und Koordinatenfunktionen \mathlk{x_1 , \ldots , x_n}{}} {} {}
in der lokalen Darstellung der Differentialform

\mavergleichskettedisp
{\vergleichskette
{ \alpha_* \omega
}
{ =} { f dx_1 \wedge \ldots \wedge dx_n
}
{ } {
}
{ } {
}
{ } {
}
}
{}{}{}
die Funktion $f$ überall positiv ist.
}{Eine
\definitionsverweis {differenzierbare}{}{}
\definitionsverweis {Differentialform}{}{}
$\omega$ auf $M$ heißt geschlossen, wenn ihre
\definitionsverweis {äußere Ableitung}{}{}

\mathl{d \omega=0}{} ist.
}{Man nennt

\mathl{\pi_{\text{vert} }   \circ TT(\gamma)}{} die
die
tangentiale Beschleunigung
von $\gamma$, wobei die Projektion vom
\definitionsverweis {Levi-Civita-Zusammenhang}{}{}
herrührt.
}

 }

__NOEDITSECTION__

\inputaufgabeklausurloesung
{3}

{

Formuliere die folgenden Sätze.
\aufzaehlungdrei{Der
\stichwort {Satz über die Normalkrümmung} {.}}{Die Formel für die Berechnung des kanonischen Volumens auf einer riemannschen Mannigfaltigkeit in einer Karte.}{Der \stichwort {Satz von Stokes} {} für Mannigfaltigkeiten mit Rand.}

}
{

\aufzaehlungdrei{\faktsituation {Es sei

\mavergleichskette
{\vergleichskette
{ W
}
{ \subseteq }{ \R^n
}
{  }{
}
{    }{
}
{   }{
}
}
{}{}{}
\definitionsverweis {offen}{}{,}
\maabb {h} { W } { \R
} {}
eine zweimal
\definitionsverweis {stetig differenzierbare Funktion}{}{}
und

\mavergleichskette
{\vergleichskette
{ Y
}
{ = }{ h^{-1}(c)
}
{  }{
}
{    }{
}
{   }{
}
}
{}{}{}
die
\definitionsverweis {Faser}{}{}
zu

\mavergleichskette
{\vergleichskette
{ c
}
{ \in }{ \R
}
{  }{
}
{    }{
}
{   }{
}
}
{}{}{,}
wobei $h$ in jedem Punkt von $Y$
\definitionsverweis {regulär}{}{}
sei. Es sei

\mavergleichskette
{\vergleichskette
{ P
}
{ \in }{ Y
}
{  }{
}
{    }{
}
{   }{
}
}
{}{}{.}
Es sei

\mavergleichskette
{\vergleichskette
{ E
}
{ \subseteq }{ \R^n
}
{  }{
}
{    }{
}
{   }{
}
}
{}{}{}
eine
\definitionsverweis {Normalebene}{}{}
durch $P$ an $Y$.}
\faktfolgerung {Dann ist die
\definitionsverweis {Normalkrümmung}{}{}
von $Y$ in $P$ gleich der
\definitionsverweis {Krümmung}{}{}
der ebenen Kurve

\mavergleichskette
{\vergleichskette
{ Y \cap (P+E)
}
{ \subseteq }{ P+E
}
{ \cong }{ \R^2
}
{    }{
}
{   }{
}
}
{}{}{}
im Punkt $P$.}
\faktzusatz {}
\faktzusatz {}}{Es sei $M$ eine
\definitionsverweis {orientierte}{}{}
\definitionsverweis {riemannsche Mannigfaltigkeit}{}{}
und $\omega$ die
\definitionsverweis {kanonische Volumenform}{}{.}
Es sei
\maabbdisp {\alpha} { U } { V
} {}
eine
\definitionsverweis {orientierte Karte}{}{}
mit

\mavergleichskettedisp
{\vergleichskette
{ V
}
{ \subseteq} { \R^n
}
{ } {
}
{ } {
}
{ } {
}
}
{}{}{}
offen mit Koordinaten
\mathl{x_1 , \ldots , x_n}{} mit der
\definitionsverweis {metrischen Fundamentalmatrix}{}{}

\mathl{G=(g_{ij})_{1 \leq i,j \leq n}}{} und
\mathl{g= \det  G}{.} Dann ist

\mavergleichskettedisp
{\vergleichskette
{  \alpha_* \omega
}
{ =} { \sqrt{ g} dx_1 \wedge \ldots \wedge dx_n
}
{ } {
}
{ } {
}
{ } {
}
}
{}{}{.}
Für eine
\definitionsverweis {messbare Teilmenge}{}{}

\mavergleichskette
{\vergleichskette
{ T
}
{ \subseteq }{ U
}
{  }{
}
{    }{
}
{   }{
}
}
{}{}{}
ist

\mavergleichskettedisp
{\vergleichskette
{
\int_{  T   }  \omega
}
{ =} {
\int_{ \alpha(T)  }   \sqrt{g} dx_1 \wedge \ldots \wedge dx_n
}
{ =} { \int_{ \alpha(T)  }   \sqrt{g}    \,  d \lambda^n
}
{ } {
}
{ } {
}
}
{}{}{.}}{Es sei $M$ eine
$n$-\definitionsverweis {dimensionale}{}{}
\definitionsverweis {orientierte}{}{}
\definitionsverweis {differenzierbare Mannigfaltigkeit mit Rand}{}{}
$\partial M$ und mit
\definitionsverweis {abzählbarer Basis der Topologie}{}{,}
und es sei $\omega$ eine
\definitionsverweis {stetig differenzierbare}{}{}
$(n-1)$-\definitionsverweis {Differentialform}{}{}
mit
\definitionsverweis {kompaktem}{}{}
\definitionsverweis {Träger}{}{} auf $M$. Dann ist

\mavergleichskettedisp
{\vergleichskette
{
\int_{ M  }  d\omega
}
{ =} {
\int_{ \partial M  }  \omega
}
{ } {
}
{ } {
}
{ } {
}
}
{}{}{.}}

 }

__NOEDITSECTION__

\inputaufgabeklausurloesung
{0}

{

}
{  }

__NOEDITSECTION__

\inputaufgabeklausurloesung
{0}

{

}
{  }

__NOEDITSECTION__

\inputaufgabeklausurloesung
{4}

{

Beweise den Satz über die Beziehung zwischen Krümmung und Einheitsnormalenvektor einer bogenparametrisierten ebenen Kurve.

}
{

Die erste Aussage folgt unmittelbar aus der Definition

\mavergleichskettedisp
{\vergleichskette
{ \kappa(t)
}
{ =} { \gamma_1'(t) \gamma_2^{\prime \prime} (t) - \gamma_2'(t) \gamma_1^{\prime \prime} (t)
}
{ } {
}
{ } {
}
{ } {
}
}
{}{}{.}
Aus

\mavergleichskettedisp
{\vergleichskette
{ \left(  \gamma_1'(t)  \right)^2 +  \left(  \gamma_2' (t)  \right)^2
}
{ =} { 1
}
{ } {
}
{ } {
}
{ } {
}
}
{}{}{}
folgt

\mavergleichskettedisp
{\vergleichskette
{ \gamma_1'(t) \gamma_1^{\prime \prime}(t) + \gamma_2'(t) \gamma_2^{\prime \prime} (t)
}
{ =} { 0
}
{ } {
}
{ } {
}
{ } {
}
}
{}{}{,}
daher steht $\gamma^{\prime \prime}(t)$ senkrecht auf $\gamma'(t)$ und ist linear abhängig zu
\mathl{N(t)}{.} Somit ist

\mavergleichskettedisp
{\vergleichskette
{ \gamma^{\prime \prime}(t)
}
{ =} { \left\langle  \gamma^{\prime \prime}(t) ,  N(t)  \right\rangle  N(t)
}
{ =} { \kappa(t) N(t)
}
{ } {
}
{ } {
}
}
{}{}{}
und somit ist

\mavergleichskettedisp
{\vergleichskette
{ \Vert { \gamma^{\prime \prime}(t) } \Vert
}
{ =} { \betrag {  \kappa(t) }
}
{ } {
}
{ } {
}
{ } {
}
}
{}{}{.}

 }

__NOEDITSECTION__

\inputaufgabeklausurloesung
{5 (1+2+2)}

{

Wir betrachten auf der zweidimensionalen
\definitionsverweis {Einheitssphäre}{}{}
$Y$ das
\definitionsverweis {tangentiale Vektorfeld}{}{}

\mavergleichskettedisp
{\vergleichskette
{ F \begin{pmatrix}  x  \\ y \\ z   \end{pmatrix}
}
{ =} { \begin{pmatrix}  y  \\ -x\\  0   \end{pmatrix}
}
{ } {
}
{ } {
}
{ } {
}
}
{}{}{.}
\aufzaehlungdrei{Bestimme die
\definitionsverweis {Jacobi-Matrix}{}{}
zu $F$ auf dem $\R^3$.
}{Bestimme für den Punkt

\mavergleichskette
{\vergleichskette
{ P
}
{ = }{ \begin{pmatrix}  0  \\ 1 \\  0   \end{pmatrix}
}
{ \in }{ Y
}
{    }{
}
{   }{
}
}
{}{}{}
eine Matrix für die Abbildung
\maabbeledisp {} {T_PY} { T_PY
} { v } { \nabla_v F
} {,}
bezüglich einer geeigneten
\definitionsverweis {Basis}{}{.}
}{Bestimme für den Punkt

\mavergleichskette
{\vergleichskette
{ Q
}
{ = }{ \begin{pmatrix}  0  \\ 0\\  1   \end{pmatrix}
}
{ \in }{ Y
}
{    }{
}
{   }{
}
}
{}{}{}
eine Matrix für die Abbildung
\maabbeledisp {} {T_QY} { T_QY
} { v } { \nabla_v F
} {,}
bezüglich einer geeigneten
\definitionsverweis {Basis}{}{.}
}

}
{

\aufzaehlungdrei{Die Jacobi-Matrix ist

\mathdisp {\begin{pmatrix}  0   &   1  &  0  \\  -1 &  0  &  0 \\ 0  &  0 &  0  \end{pmatrix}} { . }

}{Als Normalenfeld nehmen wir $\begin{pmatrix}  x  \\ y \\ z   \end{pmatrix}$, in

\mavergleichskette
{\vergleichskette
{ P
}
{ = }{ \begin{pmatrix}  0  \\ 1 \\  0   \end{pmatrix}
}
{  }{
}
{    }{
}
{   }{
}
}
{}{}{}
bilden
\mathkor {} {\begin{pmatrix}  1  \\ 0 \\  0   \end{pmatrix}} {und} {\begin{pmatrix}  0  \\ 0\\  1   \end{pmatrix}} {}
eine Basis des Tangentialraumes. Es ist

\mavergleichskettedisp
{\vergleichskette
{ \begin{pmatrix}  0   &   1  &  0  \\  -1 &  0  &  0 \\ 0  &  0 &  0  \end{pmatrix} \begin{pmatrix}  1  \\ 0 \\  0   \end{pmatrix}
}
{ =} { \begin{pmatrix}  0  \\ -1\\  0   \end{pmatrix}
}
{ } {
}
{ } {
}
{ } {
}
}
{}{}{}
und

\mavergleichskettedisp
{\vergleichskette
{  \begin{pmatrix}  0   &   1  &  0  \\  -1 &  0  &  0 \\ 0  &  0 &  0  \end{pmatrix} \begin{pmatrix}  0  \\ 0\\  1   \end{pmatrix}
}
{ =} { \begin{pmatrix}  0  \\ 0\\  0   \end{pmatrix}
}
{ } {
}
{ } {
}
{ } {
}
}
{}{}{.}
Dabei ist
\mathl{\begin{pmatrix}  0  \\ -1\\  0   \end{pmatrix}}{} ein Vielfaches zum Normalenvektor und daher ist seine orthogonale Projektion auf den Tangentialraum ebenfalls der Nullvektor. Die Abbildung
\mathl{v \mapsto \nabla_vF}{} ist also die Nullabbildung.
}{In

\mavergleichskette
{\vergleichskette
{ Q
}
{ = }{ \begin{pmatrix}  0  \\ 0\\  1   \end{pmatrix}
}
{  }{
}
{    }{
}
{   }{
}
}
{}{}{}
bilden
\mathkor {} {\begin{pmatrix}  1  \\ 0 \\  0   \end{pmatrix}} {und} {\begin{pmatrix}  0  \\ 1 \\  0   \end{pmatrix}} {}
eine Basis des Tangentialraumes. Es ist

\mavergleichskettedisp
{\vergleichskette
{ \begin{pmatrix}  0   &   1  &  0  \\  -1 &  0  &  0 \\ 0  &  0 &  0  \end{pmatrix} \begin{pmatrix}  1  \\ 0 \\  0   \end{pmatrix}
}
{ =} { \begin{pmatrix}  0  \\ -1\\  0   \end{pmatrix}
}
{ } {
}
{ } {
}
{ } {
}
}
{}{}{}
und

\mavergleichskettedisp
{\vergleichskette
{ \begin{pmatrix}  0   &   1  &  0  \\  -1 &  0  &  0 \\ 0  &  0 &  0  \end{pmatrix} \begin{pmatrix}  0  \\ 1 \\  0   \end{pmatrix}
}
{ =} { \begin{pmatrix}  1  \\ 0 \\  0   \end{pmatrix}
}
{ } {
}
{ } {
}
{ } {
}
}
{}{}{.}
Beide Bildvektoren sind bereits tangential zu $Y$ in $Q$. Eine beschreibende Matrix zu
\mathl{v \mapsto \nabla_vF}{} ist

\mathdisp {\begin{pmatrix}  0   &   -1  \\  1  &  0 \end{pmatrix}} { . }

}

 }

__NOEDITSECTION__

\inputaufgabeklausurloesung
{4 (2+2)}

{

Wir betrachten die Ellipse

\mavergleichskettedisp
{\vergleichskette
{ E
}
{ =} { \left\{ (u,v)  \mid   2u^2 +5v^2 = 4         \right\}
}
{ } {
}
{ } {
}
{ } {
}
}
{}{}{.}
\aufzaehlungdrei{Definiere eine surjektive stetig differenzierbare Abbildung
\maabb {} {\R} {E
} {.}
}{Beschreibe einen
\definitionsverweis {Diffeomorphismus}{}{}
zwischen der Sphäre $S^1$ und  $E$.
}{
}

}
{

Wir arbeiten mit der Beschreibung

\mavergleichskettedisp
{\vergleichskette
{ E
}
{ =} { \left\{ (u,v)  \mid   { \frac{   1 }{  2 } } u^2+ { \frac{   5 }{  4 } } v^2 = 1         \right\}
}
{ } {
}
{ } {
}
{ } {
}
}
{}{}{.}
\aufzaehlungzwei {Die Abbildung
\maabbeledisp {} {\R} { E
} { \alpha} { \left(  \sqrt{2} \cos  \alpha   , \,   { \frac{   2 }{  \sqrt{5}  } } \sin  \alpha                   \right)
} {}
ist als Abbildung in den $\R^2$ stetig differenzierbar und das Bild liegt auf $E$. Daher ist sie auch als Abbildung nach $E$ stetig differenzierbar. Die Surjektivität wird in (2) mitbewiesen.
} {Wir betrachten die Abbildung
\maabbeledisp {} {S^1} {E
} {(x,y)} { \left(  \sqrt{2} x  , \,   { \frac{   2 }{  \sqrt{5}  } } y                   \right)
} {.}
Diese ist als Einschränkung einer bijektiven linearen Abbildung des $\R^2$ stetig differenzierbar und injektiv und landet
\zusatzklammer {aufgrund der expliziten Gleichungen} {} {}
in der Tat in $E$. Sie ist auch surjektiv, da man direkt die Umkehrabbildung
\maabbeledisp {} { E } {S^1
} {(u,v)} { \left(  { \frac{   1 }{  \sqrt{2}  } } u  , \,   { \frac{   \sqrt{5}  }{  2  } } v                   \right)
} {,}
angeben kann. Da die trigonometrische Parametrisierung des Einheitskreisen surjektiv ist, folgt, dass auch die in (1) angegebene Abbildung surjektiv ist.
}

 }

__NOEDITSECTION__

\inputaufgabeklausurloesung
{0}

{

}
{  }

__NOEDITSECTION__

\inputaufgabeklausurloesung
{4}

{

Man gebe ein
\definitionsverweis {stetiges}{}{} \definitionsverweis {Vektorfeld}{}{}
auf $S^2$ an, das nur eine Nullstelle besitzt.

}
{

Wir betrachten auf dem $\R^2$ das stetige Vektorfeld $F$, das durch

\mavergleichskettedisp
{\vergleichskette
{F(x,y)
}
{ =} { { \frac{   1 }{  x^2+y^2 } }  e_1
}
{ } {
}
{ } {
}
{ } {
}
}
{}{}{}
gegeben sei. Dieses hat keine Nullstelle und ist stetig. Wir transportieren dieses Vektorfeld mittels der stereographischen Projektion auf

\mavergleichskette
{\vergleichskette
{  \R^2
}
{ \cong }{ S^2 \setminus \{N\}
}
{  }{
}
{    }{
}
{   }{
}
}
{}{}{}
und ergänzen es im Nordpol durch den Wert

\mavergleichskette
{\vergleichskette
{ 0
}
{ \in }{  T_NS^2
}
{  }{
}
{    }{
}
{   }{
}
}
{}{}{.}
Wir behaupten, dass dieses Vektorfeld stetig ist. Dazu sei
\mathl{P_n}{} sei eine Folge auf $S^2$, die gegen $N$ konvergiert. Dabei können wir direkt annehmen, dass alle

\mavergleichskette
{\vergleichskette
{P_n
}
{ \neq }{N
}
{  }{
}
{    }{
}
{   }{
}
}
{}{}{}
sind. Das Kartenbild dieser Folge ist

\mavergleichskettedisp
{\vergleichskette
{P'_n
}
{ =} { (x_n,y_n)
}
{ } {
}
{ } {
}
{ } {
}
}
{}{}{,}
und da $P_n$ gegen den Nordpol konvergiert, divergiert
\mathl{\betrag { P_n' }}{} bestimmt gegen $\infty$. Daher konvergiert die Folge

\mavergleichskettedisp
{\vergleichskette
{ F(P'_n)
}
{ =} { F(x_n , y_n)
}
{ =} { { \frac{   1 }{  x^2+y^2 } } e_1
}
{ } {
}
{ } {
}
}
{}{}{}
gegen $0$.

 }

__NOEDITSECTION__

\inputaufgabeklausurloesung
{6}

{

Es sei ein
\definitionsverweis {Verklebungsdatum}{}{}

\mathbed {U_i} {}
 {i \in I} {}
 {} {} {} {,}
für
\definitionsverweis {topologische Räume}{}{}
gegeben. Zeige, dass es einen eindeutig bestimmten topologischen Raum $X$, eine offene Überdeckung

\mavergleichskette
{\vergleichskette
{ X
}
{ = }{ \bigcup_{i \in I} V_i
}
{  }{
}
{    }{
}
{   }{
}
}
{}{}{}
und
\definitionsverweis {Homöomorphismen}{}{}
\maabb {\psi_i} {U_i } { V_i
} {}
derart gibt, dass

\mavergleichskettedisp
{\vergleichskette
{ \psi_i  \left(  U_{ij}  \right)
}
{ =} { V_i \cap V_j
}
{ } {
}
{ } {
}
{ } {
}
}
{}{}{}
ist und

\mavergleichskettedisp
{\vergleichskette
{ \psi_i \vert_{U_{ij} }
}
{ =} { \psi_j \vert_{U_{ji} }  \circ   \varphi_{ji}
}
{ } {
}
{ } {
}
{ } {
}
}
{}{}{}
gilt.

}
{

Es sei $Y$ die disjunkte Vereinigung der $U_i$. Wir definieren auf $Y$ eine
\definitionsverweis {Äquivalenzrelation}{}{}
$\sim$, wobei wir Punkte

\mavergleichskette
{\vergleichskette
{ x_i
}
{ \in }{ U_i
}
{  }{
}
{    }{
}
{   }{
}
}
{}{}{}
und

\mavergleichskette
{\vergleichskette
{ x_j
}
{ \in }{ U_j
}
{  }{
}
{    }{
}
{   }{
}
}
{}{}{}
als äquivalent ansehen, wenn

\mavergleichskette
{\vergleichskette
{ x_i
}
{ \in }{ U_{ij}
}
{  }{
}
{    }{
}
{   }{
}
}
{}{}{,}

\mavergleichskette
{\vergleichskette
{ x_j
}
{ \in }{ U_{ji}
}
{  }{
}
{    }{
}
{   }{
}
}
{}{}{}
und

\mavergleichskette
{\vergleichskette
{ \varphi_{ji} (x_i)
}
{ = }{ x_j
}
{  }{
}
{    }{
}
{   }{
}
}
{}{}{}
ist. Die Eigenschaften einer Äquivalenzrelation sind dabei durch die Kozykelbedingung gesichert, siehe
[[Topologischer Raum/Verklebungsdatum/Äquivalenzrelation/Aufgabe|[[:Kurs:Bündel, Garben und Kohomologie (Osnabrück 2019-2020)/4/Klausur mit Lösungen/latex]] (Bündel, Garben und Kohomologie (Osnabrück 2019-2020))[[Kategorie:Kurs:Differentialgeometrie/Referenznummer|Bündel, Garben und Kohomologie (Osnabrück 2019-2020)]] (Differentialgeometrie (Osnabrück 2023))]].
Wir setzen

\mavergleichskettedisp
{\vergleichskette
{X
}
{ \defeq} { Y/ \sim
}
{ } {
}
{ } {
}
{ } {
}
}
{}{}{}
und versehen $X$ mit der
\definitionsverweis {Quotiententopologie}{}{.}
Die Verknüpfungen

\mathdisp {U_i \longrightarrow Y \longrightarrow X} {  }

sind die $\psi_i$, und $V_i$ sind die Bilder dieser Abbildungen. Daher liegen Homöomorphismen
\maabb {\psi_i} {U_i } { V_i
} {}
vor. Dabei ist zu

\mavergleichskette
{\vergleichskette
{ x
}
{ \in }{ U_i
}
{  }{
}
{    }{
}
{   }{
}
}
{}{}{}

\mavergleichskettedisp
{\vergleichskette
{ \psi_i (x)
}
{ \in} { V_j
}
{ } {
}
{ } {
}
{ } {
}
}
{}{}{}
genau dann, wenn

\mavergleichskette
{\vergleichskette
{ x
}
{ \in }{ U_{ij}
}
{  }{
}
{    }{
}
{   }{
}
}
{}{}{}
ist, da genau in diesem Fall $x$ mit
\mathl{\varphi_{ij}(x)}{} identifiziert wird. Daher ist

\mavergleichskette
{\vergleichskette
{ \psi_i(U_{ij})
}
{ = }{ V_i \cap V_j
}
{  }{
}
{    }{
}
{   }{
}
}
{}{}{.}
Die Kommutativität des Diagramms

\mathdisp {\begin{matrix} U_{ij}   & \stackrel{ \varphi_{ji} }{\longrightarrow} &  U_{ji}   & \\ & \!\!\! \!\! \psi_i \searrow & \downarrow \psi_j \!\!\! \!\! & \\ & &  V_i \cap V_j   & \!\!\!\!\! \!\!\!  \\ \end{matrix}} {  }

folgt ebenso.

 }

__NOEDITSECTION__

\inputaufgabeklausurloesung
{2}

{

Es sei $M$ eine
\definitionsverweis {differenzierbare Mannigfaltigkeit}{}{,}
\maabb {\gamma} { I } { M
} {}
eine
\definitionsverweis {differenzierbare Kurve}{}{}
in $M$ und $\omega$ eine differenzierbare
$1$-\definitionsverweis {Form}{}{}
auf $M$. Zeige, dass

\mavergleichskettedisp
{\vergleichskette
{ \gamma^* (d \omega)
}
{ =} { 0
}
{ } {
}
{ } {
}
{ } {
}
}
{}{}{}
ist.

}
{

Es ist $d \omega$ eine $2$-Form auf $M$  und das überträgt sich auf $\gamma^*(d \omega)$. Auf einer eindimensionalen Mannigfaltigkeit sind aber  alle $2$-Formen gleich $0$.

 }

__NOEDITSECTION__

\inputaufgabeklausurloesung
{6 (3+3)}

{

Es sei $M$ eine
\definitionsverweis {kompakte}{}{}
\definitionsverweis {differenzierbare Mannigfaltigkeit}{}{}
mit zumindest zwei Punkten.
\aufzaehlungzwei {Zeige, dass eine
\definitionsverweis {exakte}{}{}
stetige
$1$-\definitionsverweis {Differentialform}{}{}
auf $M$ zumindest zwei Nullstellen besitzt.
} {Zeige, dass dies für eine
\definitionsverweis {geschlossene}{}{}
Differentialform nicht gelten muss.
}

}
{

\aufzaehlungzwei {Eine exakte Differentialform kann man als

\mavergleichskettedisp
{\vergleichskette
{ \omega
}
{ =} { df
}
{ } {
}
{ } {
}
{ } {
}
}
{}{}{}
mit einer differenzierbaren Funktion
\maabbdisp {f} { M } {\R
} {}
ansetzen. Wenn $f$ konstant ist, ist

\mavergleichskette
{\vergleichskette
{ df
}
{ = }{ 0
}
{  }{
}
{    }{
}
{   }{
}
}
{}{}{}
und die Aussage ist klar. Sei $f$ also nicht konstant. Nach
[[Kompaktheit/Stetige Funktion/Maximum wird angenommen/Fakt|Fakt]] [[Kurs:Differentialgeometrie/Kompaktheit/Stetige Funktion/Maximum wird angenommen/Fakt/Faktreferenznummer|*****]]
nimmt die stetige Funktion ein globales Maximum und ein globales Minimum an, sagen wir in den Punkten
\mathkor {} {P} {bzw.} {Q} {.}
Nach
[[Differenzierbare Mannigfaltigkeit/Funktion/Extrema/Differentialform/Nullstelle/Aufgabe|Aufgabe]] [[Kurs:Differentialgeometrie/Differenzierbare Mannigfaltigkeit/Funktion/Extrema/Differentialform/Nullstelle/Aufgabe/Aufgabereferenznummer|*****]]
ist dann

\mavergleichskettedisp
{\vergleichskette
{ (df)_P
}
{ =} { (df)_Q
}
{ =} { 0
}
{ } {
}
{ } {
}
}
{}{}{.}
} {Wir betrachten auf der kompakten eindimensionalen Sphäre $S^1$ die Differentialform

\mavergleichskettedisp
{\vergleichskette
{ \omega(x,y)
}
{ =} {  ydx - xdy
}
{ } {
}
{ } {
}
{ } {
}
}
{}{}{.}
Dies ist in der Tat eine Differentialform auf dem Kreis, da
\mathl{\left(  y   , \,   -x                   \right)}{} tangential zum Ortsvektor
\mathl{(x,y)}{} ist. Diese Form ist geschlossen, da wir in der Dimension $1$ sind. Sie hat keine Nullstelle, da auf dem Kreis nicht
\mathkor {} {x} {und} {y} {}
gleichzeitig $0$ sind.
}

 }

__NOEDITSECTION__

\inputaufgabeklausurloesung
{4}

{

Man gebe ein Beispiel für eine
\definitionsverweis {zusammenhängende}{}{}
\definitionsverweis {Mannigfaltigkeit mit Rand}{}{,}
deren Rand aus
\zusatzklammer {einer disjunkten Vereinigung von} {} {}
unendlich vielen  $\R^1$ besteht.

}
{

Die Funktion
\maabbeledisp {h} {]-1,1[ } { \R
} { x } { { \frac{   1  }{  x^2-1 } }
} {,}
läuft an den beiden Grenzen gegen unendlich, der Graph zu $h$ ist diffeomorph zu $]-1,1[$ und damit auch zu $\R$. Wir betrachten nun zu

\mavergleichskette
{\vergleichskette
{ n
}
{ \in }{  \Z
}
{  }{
}
{    }{
}
{   }{
}
}
{}{}{}
die entsprechende Funktion, deren Definitonsbereich um $4n$ verschoben wird
\zusatzklammer {also mit dem Definitionsbereich
\mathl{]4n-1,4n+1[}{}} {} {.}
Es sei $G_n$ der zugehörige Graph und $U_n$ der zugehörige offene
\zusatzklammer {!} {} {}
Epigraph. Wir setzen

\mavergleichskettedisp
{\vergleichskette
{ M
}
{ =} { \R^2 \setminus \biguplus_{n \in \Z} U_n
}
{ } {
}
{ } {
}
{ } {
}
}
{}{}{.}
Diese ist eine zusammenhängende Mannigfaltigkeit mit Rand und der Rand ist die disjunkte Vereinigung der $G_n$.

 }

__NOEDITSECTION__

\inputaufgabeklausurloesung
{10}

{

Beweise den  \stichwort {Satz über die Partition der Eins} {.}

}
{

Nach
[[Mannigfaltigkeit/Abzählbare Basis/Lokal endlicher Atlas/Fakt|Lemma 22.9 (Differentialgeometrie (Osnabrück 2023))]]
können wir davon ausgehen, dass eine offene Überdeckung aus Kartengebieten

\mathbed {U_j} {}
 {j \in J} {}
 {} {} {} {,}
\zusatzklammer {$J$ abzählbar} {} {}
mit
\maabbdisp {\alpha_j} {U_j } { V_j
} {}
und mit Ballumgebungen

\mavergleichskettedisp
{\vergleichskette
{ U \left(   0 , \delta_j   \right)
}
{ \subset} { B  \left(  0 , \epsilon_j  \right)
}
{ \subset} { V_j
}
{ } {
}
{ } {
}
}
{}{}{}
\zusatzklammer {mit

\mavergleichskettek
{\vergleichskettek
{ \delta_j
}
{ < }{ \epsilon_j
}
{  }{
}
{    }{
}
{   }{
}
}
{}{}{}} {} {}
vorliegt derart, dass auch die
\mathl{\alpha_j^{-1} \left(  U \left(   0 , \delta_j   \right)  \right)}{} eine Überdeckung von $M$ bilden und dass jeder Punkt

\mavergleichskette
{\vergleichskette
{ P
}
{ \in }{ M
}
{  }{
}
{    }{
}
{   }{
}
}
{}{}{}
nur in endlich vielen der $U_j$ und insbesondere nur in endlich vielen dieser
\mathl{\alpha_j^{-1} \left(  U \left(   0 , \delta_j   \right)  \right)}{} enthalten ist. Auf $V_j$ betrachten wir die Funktion $g_j$, die durch

\mavergleichskettedisp
{\vergleichskette
{ g_j (v)
}
{ =} { \begin{cases}  e^{ - { \frac{   1  }{  \left(  \delta_j^2- \Vert { v } \Vert^2  \right)^2 } } }  \text{ für } \Vert { v } \Vert < \delta_j \, , \\ 0 \text{ sonst} \, , \end{cases}
}
{ } {
}
{ } {
}
{ } {
}
}
{}{}{}
definiert ist. Diese Funktion hat genau auf
\mathl{U \left(   0 , \delta_j   \right)}{} einen positiven Wert und ihr
\definitionsverweis {Träger}{}{}
ist
\mathl{B  \left(  0 , \delta_j  \right)}{.} Eine Betrachtung auf den beiden offenen Teilmengen
\zusatzklammer {die $V_j$ überdecken} {} {}
\mathkor {} {U \left(   0 , \epsilon_j   \right)} {und} {V_j \setminus B  \left(  0 , \delta_j  \right)} {}
zeigt, dass $g_j$ unendlich oft differenzierbar ist. Wir definieren eine Funktion
\maabbdisp {\tilde{g}_j} { M } {\R
} {}
durch

\mavergleichskettedisp
{\vergleichskette
{ \tilde{g}_j(x)
}
{ =} { \begin{cases}  g( \alpha_j(x)) \text{ für }  x \in \alpha_j^{-1}( U \left(   0 , \delta_j   \right) ) \, , \\  0 \text{ sonst} \, . \end{cases}
}
{ } {
}
{ } {
}
{ } {
}
}
{}{}{}
Diese Funktion ist
\definitionsverweis {stetig differenzierbar}{}{}
auf $M$, da der \anfuehrung{Streifen}{}
\mathl{B  \left(  0 , \epsilon_j  \right) \setminus U \left(   0 , \delta_j   \right)}{} einen glatten Übergang erlaubt. Wir setzen

\mavergleichskettedisp
{\vergleichskette
{ \tilde{g}(x)
}
{ \defeq} { \sum_{j \in J} \tilde{g}_j(x)
}
{ } {
}
{ } {
}
{ } {
}
}
{}{}{,}
wobei dies für jeden Punkt eine endliche Summe ist, da der
\definitionsverweis {Träger}{}{}
von $\tilde{g}_j$ in

\mavergleichskettedisp
{\vergleichskette
{ \alpha^{-1} \left(  B  \left(  0 , \delta_j  \right)  \right)
}
{ \subset} { U_j
}
{ } {
}
{ } {
}
{ } {
}
}
{}{}{}
liegt. Diese Funktion ist
\definitionsverweis {stetig differenzierbar}{}{}
auf $M$ und überall positiv, da die
\mathl{\tilde{g}_j(x)}{} auf den überdeckenden Mengen
\mathl{\alpha^{-1} \left(  U \left(   0 , \delta_j   \right)  \right)}{} positiv sind. Dann bilden die

\mavergleichskettedisp
{\vergleichskette
{ h_j
}
{ =} { { \frac{   \tilde{g}_j  }{ \tilde{g}  } }
}
{ } {
}
{ } {
}
{ } {
}
}
{}{}{}
die gesuchte Partition der Eins.

 }

__NOEDITSECTION__

\inputaufgabeklausurloesung
{10 (2+6+2)}

{

Wir betrachten die
$2$-\definitionsverweis {Differentialform}{}{}

\mavergleichskettedisp
{\vergleichskette
{ \omega
}
{ =} { x dy \wedge dz -y dx \wedge dz + z dx \wedge dy
}
{ } {
}
{ } {
}
{ } {
}
}
{}{}{}
auf der Einheitskugel
\mathl{B  \left(  0 , 1  \right)}{.}
\aufzaehlungdreiabc{Zeige, dass $d \omega$ das Dreifache der Standardvolumenform auf der Kugel ist.
}{Zeige, dass
\mathl{\omega\vert_{S^2}}{} die Standardflächenform auf der Einheitssphäre ist.
}{Berechne die Kugeloberfläche aus dem
[[Allgemeines Kugelvolumen/Mit Cavalieri-Prinzip/Beispiel|Kugelvolumen]]
mit dem
[[Mannigfaltigkeit mit Rand/Satz von Stokes/Fakt|Satz von Stokes]].
}

}
{

a) Es ist

\mavergleichskettealign
{\vergleichskettealign
{d \omega
}
{ =} { d \left(  x dy \wedge dz  -y dx \wedge dz +  z dx \wedge dy \right)
}
{ =} { dx \wedge dy \wedge dz - dy \wedge dx \wedge dz + d z \wedge dx \wedge dy
}
{ =} { dx \wedge dy \wedge dz + dx \wedge dy \wedge dz + d x \wedge dy \wedge dz
}
{ =} { 3 d x \wedge dy \wedge dz
}
}
{}
{}{,}
wobei wir verwendet haben, dass sich das Vorzeichen bei der Vertauschung von Faktoren im Dachprodukt ändert.

b) Für einen Punkt
\mathl{P = \begin{pmatrix}  x  \\ y \\  z   \end{pmatrix}  \in S^2}{} und zwei orthonormale,
\zusatzklammer {zusammen mit $P$} {} {}
die Orientierung repräsentierende Tangentialvektoren
\mathl{u= \begin{pmatrix} u_1  \\u_2 \\ u_3   \end{pmatrix}}{} und
\mathl{v= \begin{pmatrix} v_1  \\v_2 \\ v_3   \end{pmatrix}}{} ist nach
[[Dachprodukt/Natürliche Dualität/Endlichdimensional/Fakt|Satz 58.4 (Lineare Algebra (Osnabrück 2024-2025))]]

\mavergleichskettealign
{\vergleichskettealign
{ \omega \vert_{S^2}  (u,v)
}
{ =} { \omega (u,v)
}
{ =} { z \det  \begin{pmatrix} u_1   &  v_1   \\ u_2  & v_2  \end{pmatrix} - y \det  \begin{pmatrix} u_1   &  v_1   \\ u_3  & v_3  \end{pmatrix} + x \det  \begin{pmatrix} u_2   &  v_2   \\ u_3  & v_3  \end{pmatrix}
}
{ =} { \det  \begin{pmatrix}  x   &  u_1  & v_1  \\  y  & u_2  & v_2  \\ z  & u_3  & v_3  \end{pmatrix}
}
{ } {
}
}
{}
{}{.}
Die Determinante eines die Orientierung repräsentierenden Orthonormalsystems ist aber $1$. Also erfüllt die eingeschränkte Differentialform die Eigenschaft, die die
\definitionsverweis {Standardvolumenform}{}{}
charakterisiert.

c) Die Oberfläche der Kugel ist unter Verwendung von a), b), des Satzes von Stokes und des Kugelvolumens gleich

\mavergleichskettedisp
{\vergleichskette
{ \int_{S^2} \omega
}
{ =} { \int_{ B  \left(  0, 1 \right)  }  d \omega
}
{ =} { \int_{ B  \left(  0, 1 \right)  } 3 d x \wedge dy \wedge dz
}
{ =} { 3 \lambda^3 (   B  \left(  0, 1 \right)     )
}
{ =} { 4 \pi
}
}
{}{}{.}

 }

__NOEDITSECTION__

\inputaufgabeklausurloesung
{0}

{

}
{  }

 [[Kategorie:Latexseite]]
